% =====================================================================
%  ME 697 -- bitwhiz: Reactive Burn, the Lead Shock                 v1
%  v1 (Nov 7, 2026, with Claude): the front matter and the first two
%  boxes, for DK's review, beside ThermalCascade_rev-0 (v2).  DK's
%  decisions (Nov 7): its own sheet, shared by the burn models that use
%  the lead shock (ThermalCascade now; FlashFire and StochasticIgnite to
%  come); student-facing, not part of the class project; Menikoff and
%  Shaw's h written as calligraphic H (h is U3's areal mass coordinate).
%
%  VERIFIED: Eqs. (10) and (11) and Sec. 4's two cases (zero crossing,
%  local minimum) against the rendered pages 5-6 of Menikoff and Shaw
%  (2010); its DOI against page 1.  H = q (V0 - V)/2 on a steady AV
%  profile: derived (steady mass, momentum, energy) and checked to
%  round-off by bitwhiz's tests/test_lshk.f95, check 1.
%
%  STUB CONVENTION: \tbd{...} RED = still DK's decision.  Black =
%  verified.
%
%  PLANNED (not drafted): finding the lead shock on a mesh (armed above
%  a floor; marked at a zero crossing, a local minimum or a stall; P_S
%  frozen), with the measured reasons for each rule (bitwhiz's
%  doc/design/burn_models.md, "As built: step 4"); what to expect
%  (P_S against exact Hugoniot states, a smooth ramp's foot); a figure
%  of H(t) for one cell, drawn from a solved run; references.
% =====================================================================
\documentclass[11pt,addpoints]{exam}
\usepackage{amsfonts,amssymb,amsmath, amsthm}
\usepackage{graphicx}
\usepackage{enumitem}
\usepackage{systeme}
\usepackage{pgf,tikz,pgfplots}
\pgfplotsset{compat=1.15}
\usepgfplotslibrary{fillbetween}
\usepackage{mathrsfs}
\usepackage{multicol}
\usepackage{booktabs}
\usepackage[many]{tcolorbox}
\usepackage{bm}
\usepackage{cancel}
\allowdisplaybreaks

% Good Volume Symbol
\newcommand{\volume}{{\ooalign{$V$\cr\hfil\raisebox{0.35em}{\rule{2.1mm}{0.1ex}}\hfil\cr}}}

%%% Strike-out font
\usepackage{ulem}
\normalem
\pagestyle{headandfoot}

% ---- Warm theme to match the figures (U5 onward) ----
\definecolor{slate}{HTML}{2C3E50}
\definecolor{oiblue}{HTML}{0072B2}
\tcbset{colback=oiblue!4!white, colframe=slate, boxrule=0.5pt, arc=1.5mm}

% ---- Stub marker: red = placeholder (delete when the sheet is final) ----
\newcommand{\tbd}[1]{\textcolor{red}{#1}}

\firstpageheader{ME 697 Equations}{bitwhiz: Reactive Burn, the Lead Shock}{Fall 2026 -- D. Kittell}
\runningheader{ME 697 Equations}{bitwhiz: Reactive Burn, the Lead Shock}{Fall 2026 -- D. Kittell}
\runningheadrule
\firstpagefooter{}{}{}
\runningfooter{}{}{}

\makeatletter
\renewcommand*\env@matrix[1][*\c@MaxMatrixCols c]{%
  \hskip -\arraycolsep
  \let\@ifnextchar\new@ifnextchar
  \array{#1}}
\makeatother



\begin{document}

\noindent
\bf{Instructor's Note:}\normalfont\\

\noindent
\it{Not part of the class project, and not extra credit. Three of the burn models in bitwhiz, the course's open-source hydrocode, need the pressure of the first shock to reach each cell, and burn only after it arrives. A hydrocode smears every shock over a few cells, so which state is ``the shock''? Menikoff and Shaw's answer uses nothing but the cell's own history and the Hugoniot jump condition for energy \cite{menikoff2010}. This sheet gives it, with the rules bitwhiz adds to make it work on a mesh.}\normalfont\vspace{1.5em}

\textbf{These equations build upon U5. Rankine-Hugoniot Analysis and the Class Project: Hydrocode sheets (Method AV and Method RS)}\vspace{1em}


\begin{tcolorbox}[breakable]
\textbf{What the Lead Shock Adds to Your Hydrocode:}\vspace{.5em}
\begin{itemize}
\item Two numbers per cell, \hspace{3mm} \bm{$t_S$}, the time the lead shock arrives, \hspace{2mm} and \bm{$P_S$}, its pressure: \hspace{2mm} found from the cell's own history, then frozen\vspace{.5em}
\item A gate, \hspace{3mm} a burn model that uses \bm{$P_S$} burns nothing before \bm{$t_S$} \hspace{7mm} ThermalCascade, FlashFire, StochasticIgnite\vspace{.5em}
\item Frozen means frozen: \hspace{2mm} a second shock does not change \bm{$P_S$}; it acts only through the cell's own \bm{$P$} and \bm{$T$} \hspace{3mm} (one route to shock desensitization)\vspace{.5em}
\item Carried from step to step, \hspace{3mm} five numbers per cell: \hspace{2mm} \bm{$\mathcal{H}$} at the last step, its largest value, a step flag, \bm{$P_S$}, \bm{$t_S$}\vspace{.5em}
\item Notation, \hspace{3mm} Hugoniot function \bm{$\mathcal{H}$} (Menikoff and Shaw's \bm{$h$}; not U3's areal mass coordinate \bm{$h$}); \hspace{2mm} the state ahead of the shock \bm{$\left(\;\;\right)_0$}; \hspace{2mm} artificial viscosity \bm{$q$} (Method AV)\vspace{.5em}
\item \textbf{\emph{One check per cell per cycle, after the pressures: the hydro itself is unchanged}}
\end{itemize}
\end{tcolorbox}


\clearpage

\begin{tcolorbox}[breakable]
\textbf{The Hugoniot Function \cite{menikoff2010}:}\vspace{.5em}
\begin{itemize}
\item Energy jump, \hspace{3mm} \bm{$\displaystyle e_H-e_0\;=\;\frac{1}{2}\left(P_H+P_0\right)\left(\text{V}_0-\text{V}_H\right)$} \hspace{5mm} U5, any shock from the state \bm{$\left(\;\;\right)_0$}\vspace{1.0em}
\item Hugoniot function, \hspace{3mm} \bm{$\displaystyle \mathcal{H}\left(\text{V},e\right)\;=\;e-e_0-\frac{1}{2}\left(\text{V}_0-\text{V}\right)\left[P\left(\text{V},e\right)+P_0\right]$} \hspace{5mm} zero on the Hugoniot\vspace{1.0em}
\item Erratum, \hspace{3mm} Menikoff and Shaw print \bm{$P-P_0$} in their Eq.~(10); \hspace{2mm} the jump condition has \bm{$P+P_0$}, \hspace{2mm} and at \bm{$P_0=0$} the two agree\vspace{1.0em}
\item At rest, \hspace{3mm} every cell starts at \bm{$\rho_0$} with \bm{$e_0=0$}, \bm{$P_0=0$}: \hspace{5mm} \bm{$\displaystyle \mathcal{H}\;=\;e-\frac{1}{2}\left(\text{V}_0-\text{V}\right)P$}\vspace{1.0em}
\item Each cell, after each step, \hspace{3mm} \bm{$\mathcal{H}^{n+1}=\mathcal{H}\left(\text{V}^{n+1},e^{n+1}\right)$} \hspace{5mm} their Eq.~(11): a history, cell by cell\vspace{1.0em}
\item Inside a steady smeared shock, \hspace{3mm} \bm{$\sigma=P+q$} is linear in \bm{$\text{V}$} (the Rayleigh line), and steady energy gives \bm{$\displaystyle e-e_0=\frac{1}{2}\left(\sigma+P_0\right)\left(\text{V}_0-\text{V}\right)$}, so\vspace{.5em} \\
.~~~~~~~~~~\bm{$\displaystyle \mathcal{H}\;=\;\frac{1}{2}\,q\left(\text{V}_0-\text{V}\right)\;\ge\;0$} \hspace{5mm} it rises through the shock and returns to zero behind it\vspace{1.0em}
\item Smooth, isentropic compression from rest, \hspace{3mm} \bm{$\mathcal{H}\le 0$}: \hspace{2mm} the chord from \bm{$\left(\text{V}_0,0\right)$} to \bm{$\left(\text{V},P\right)$} lies above the convex isentrope beneath it\vspace{.5em}
\item \textbf{\emph{From rest, a cell's \bm{$\mathcal{H}$} first turns positive in a shock: the lead shock is where it first rises and falls back}} \hspace{3mm} (a release behind a shock keeps \bm{$\mathcal{H}>0$}: the shock's entropy stays)
\end{itemize}
\end{tcolorbox}


\clearpage

\begin{tcolorbox}[breakable]
\textbf{References:}\vspace{.5em}

% Ledger: README.md in this folder.  One \bibitem per ledger row, same key.
\begingroup
\renewcommand{\section}[2]{}% suppress the "References" heading inside the box
\begin{thebibliography}{9}
\setlength{\itemsep}{0.25em}
\bibitem{menikoff2010}
R.~Menikoff and M.~S. Shaw, ``Reactive burn models and ignition \& growth concept,'' \emph{EPJ Web of Conferences} \textbf{10}, 00003 (2010), doi:10.1051/epjconf/20101000003.
\end{thebibliography}
\endgroup
\end{tcolorbox}

\end{document}
